购买系统中，Alice重新打开订单页面时，入口先取得P1的查询投影，再显示付款与履约进度。查询投影可以由Redis承载，原购买事实继续由所属逻辑分区的完整Raft状态机保存。两条读取路径回答不同问题：页面查询P1现在显示什么状态，原操作接口根据完整身份
\code{S=(1,purchase.create.external.v2,E1,K1)}返回当初保存的输入和结果。把页面缓存全部清空以后，原C1、T1、平台键\code{pay-C1}和平台交易X1之间的关联继续存在。

下文沿用这一笔既有购买的P1/C1/T1简记。第17章为新分配对象使用带逻辑分区前缀的编号，搬迁和恢复既有对象时保留它们的原编号与原平台键。

这一职责安排决定了Redis子系统的故障处理。查询缓存丢失时，投影服务从权威事实重建。缓存暂时陈旧时，页面可以带上投影版本和更新时间，资金判断则回到权威保存域。实现细节以Redis 8.10.2为准~\cite{redis8102-source}\footnote{本章固定Redis 8.10.2、commit \code{498ecd0d6d007db11ddb3aea9428552598a78622}。Redis 8提供RSALv2、SSPLv1或AGPLv3选择；较早版本的BSD许可说明不适用于这一版本。版本、原文和源码范围保存在配套研究记录。}。

本章沿两条任务展开。第一条是购买系统采用的查询缓存路径：用回填租约阻挡失效后的旧写，持久失效事件接住发送者退出，\code{min\_revision}把调用方所要求的最低观察版本接到回源读取。RDB、AOF和复制可以缩短投影恢复时间，投影恢复点落后时由权威重建。第二条是架构选型反例：假设另一项业务把 \code{accepted:\{b1\}:S1}、\code{pending:\{b1\}:C1} 等受理与待办直接放在Redis，逐步检查快照遗漏、AOF同步、主副本接替和迁槽超时会丢掉哪些已经承担的责任。后面使用 $A_J/R_J$、C1或关闭会话推演的历史都属于这条假设路径，借它检验Redis能否满足该业务要求的持久性和接管条件。

\CEighteenFigure{ch18-01}{购买权威、原结果接口与查询缓存的读取关系}{fig:18-authority-cache}

\subsection{先认识键、命令和回复}

Redis是通过网络提供键值访问的服务器。客户端保存连接，发送命令名与参数，服务器在选定数据库的键空间中查找键并返回结果。键是用于定位数据的字节串，值带有类型：字符串可以保存一份JSON投影，哈希可以保存字段集合，列表可以保存有序元素。本节先用字符串表达P1的完整页面状态，键名\code{projection:P1}由应用约定。

下面的文字按一次请求、一次回复排列，\code{>}后是客户端命令，下一行是服务器回复。引号表示该参数包含空格或特殊字符，客户端库负责把参数编码成网络协议。首次执行时假定数据库中还没有这个键。

\begin{BookCode}
> SET projection:P1 '{"purchase_id":"P1","revision":1,"state":"paid"}' EX 60
OK
> GET projection:P1
'{"purchase_id":"P1","revision":1,"state":"paid"}'
> TTL projection:P1
(integer) 59
> DEL projection:P1
(integer) 1
> GET projection:P1
(nil)
\end{BookCode}

\code{SET}把整个字符串放到键下，同名旧字符串会被替换。\code{EX 60}同时设置60秒过期时间。\code{GET}返回当前字符串，缺键返回空值，示例用\code{(nil)}表示它。\code{TTL}返回剩余秒数，实际读数取决于经过的时间，59只是这条历史中的读数；返回$-2$表示键不存在，$-1$表示键存在但没有过期时间。\code{DEL}返回实际删除的键数，此处1说明删除了一个键。回复中的整数、空值和字符串是不同结果类型，应用不应只检查连接上是否收到了一行文字~\cite{redis8102-commands}。

\code{SET key value NX}中的\code{NX}要求键当前不存在，条件满足返回\code{OK}，条件不满足返回空值。它常被用来竞争一个缓存回填位置。服务究竟受理了哪项业务，则由该值的完整结构和恢复协议决定。只留下一个值为1的占位，无法说明它绑定了什么输入、是否已经付款以及下一步由谁恢复。

哈希适合按名称读取字段。例如\code{HSET projection:P1 revision 2 state fulfilled}在同一条命令中更新两个字段，回复是新增加的字段数，更新已有字段不增加这个数。\code{HGETALL projection:P1}返回字段与值的对应关系。键原先若保存字符串，再向它执行\code{LPUSH}或\code{HSET}，服务器会返回\code{WRONGTYPE}错误，提示当前对象类型与命令要求不符。应用选择字符串或哈希以后，要让读写方遵守同一结构。

\CEighteenFigure{ch18-19-basic}{一条命令如何定位键、检查类型并形成有类型的回复}{fig:18-basic-command}

这些人可读的示例在连接上由RESP承载。RESP是Redis的请求与回复编码协议，命令通常编码为参数数组，每个参数带长度，二进制值由长度界定。\code{OK}、整数、字符串、空值和错误有不同的协议表示，客户端库再把它们映射成相应语言对象。流水线允许客户端连续发送多条命令而不逐条等待网络往返，服务器仍分别执行和回复；流水线自身没有把多条命令合成一个事务。

\subsection{从套接字到命令执行}

客户端发送\code{GET projection:P1}，字节首先进入连接的输入缓冲区。事件循环根据可读事件调用网络处理，解析器识别请求格式，检查参数长度，并组装命令及参数。流水线中的多个完整请求可以先后解析，半个请求则留在缓冲区等待后续字节。解析完成以后，服务器已经知道调用方要求执行什么，命令随后进入执行检查。

Redis 8.10.2可以让I/O线程承担连接读写与解析相关工作。普通命令的数据访问和执行继续交回主线程，因此要分别观察解析与数据修改发生在哪个线程。主线程取得待执行请求后，检查命令、参数、权限、内存限制、复制角色和集群路由等条件，然后调用命令实现。网络处理、后台工作与命令执行形成不同路径，I/O并行没有把普通写命令变成任意并行修改同一键的\mbox{操作~\cite{redis8102-source}。}

键空间保存键到对象的关联。当前版本通过\code{kvstore}组织数据库键空间和键级过期索引，对象另带类型与编码，较短字符串可以采用紧凑编码，较大字符串使用独立的动态字符串。哈希、列表、集合等类型的内部编码服务于访问和内存成本。对于P1的投影，应用可以选择一个序列化对象，也可以选择若干哈希字段，哪个命令形成一次可见更新随之确定。一个JSON字符串能让读者一次取得整份投影，多次独立字段写则可能向读者暴露更新中的组合。

普通命令依主线程的执行次序观察和修改内存。命令产生的回复先进入输出缓冲，真正写入客户端套接字可能在稍后的事件处理或I/O阶段发生。修改还会按传播规则进入追加文件AOF与复制路径，前者供重启恢复，后者把数据变化送到副本。由此出现几个不同位置：服务器执行了更新、形成了回复、客户端收到了回复、本地同步完成、副本确认了相应复制位置。即使这些动作在正常运行中相隔很短，故障也可能落在它们之间。

查询投影可以保存为\code{\{purchase\_id:P1,revision:2,state:fulfilled\}}，一次替换整个值时，读者取得旧对象或新对象。若应用把\code{revision}和\code{state}拆成两个键，先写版本2，再写\code{fulfilled}，另一读者就可能在两次命令之间取得版本2与\code{paid}的组合。单条命令的串行执行只覆盖那条命令，数据组织直接影响接口能解释哪些观察。

同一个服务器还会处理过期删除、淘汰以及模块触发的修改。这些动作可能和原命令一起产生待传播操作，传播层因此在执行单元结束时整理修改。例如，一个键在读取时被判断为过期，传播需要反映实际删除，让AOF恢复与副本取得相应状态。维护型动作、客户端命令和业务动作在来源上不同，最终都进入本地状态变化与传播路径。

\CEighteenFigure{ch18-02}{命令执行、传播、输出与持久确认的不同位置}{fig:18-command-path}

\subsection{过期和淘汰改变的是缓存内容}

投影设置60秒TTL，表示这个缓存项在到期后不再作为有效值返回。Redis既在访问键时判断过期，也通过主动过期循环回收一部分到期数据，主动工作受时间预算约束。物理内存中暂时还有某个键，与应用可以读到有效值是两个观察。复制副本的过期删除还受主节点传播的删除和自身角色影响，主副本据此衔接过期历史。

内存达到\code{maxmemory}限制后，淘汰策略根据允许的键集合及其访问、频度或TTL等信息选择牺牲项。采用\code{noeviction}时，服务器通过拒绝相应写入应对无法释放的内存压力。缓存调用方因而要处理miss、写入失败和容量变化。对于查询投影，这些情况都可以转回投影来源；对于唯一的S绑定和原结果，任意淘汰会直接删去服务恢复旧身份所依赖的事实。

设Redis中的P1已经过期，权威中保存原输入Q7、3000分及原结果X1。调用方重发原S时，入口把原S与原输入交给所属分区。权威先找到已有绑定，返回同一个P1/C1/T1和原付款结果。调用方把金额改为3500分，却继续使用原S，则由权威报告输入冲突。这两条分支读取的是持久绑定，缓存缺失只触发页面重建。

\subsection{多命令执行与业务状态闭合}

若同一份投影需要同时修改状态和版本，可以用单命令写整份对象，也可以用事务或脚本组成服务器内的执行单元。这里的执行单元指一段连续完成的命令执行：普通命令各形成一个单元，事务则把已排队的子命令连续运行。下面假定两个键均可写，展示同一连接中的完整过程。

\begin{BookCode}
> MULTI
OK
> SET projection:P1 '{"revision":2,"state":"fulfilled"}'
QUEUED
> SET projection-version:P1 2
QUEUED
> EXEC
1) OK
2) OK
\end{BookCode}

\code{MULTI}使连接进入事务排队状态，后续普通命令回复\code{QUEUED}，说明已排队。\code{EXEC}才依次执行队列，返回按子命令顺序排列的结果数组。其他客户端的普通命令不会插到这两条子命令之间。\code{DISCARD}用于在执行前丢弃队列，回复\code{OK}。把排队回复当成写入成功，会把尚未执行的命令误报给业务调用方~\cite{redis8102-transactions}。

\code{WATCH}提供带条件的执行。例如客户端先执行\code{WATCH projection:P1}并读取旧值，再构造事务。如果受监视键在\code{EXEC}以前被修改或过期等，\code{EXEC}放弃执行队列并返回空结果。客户端重新读取和判断后才决定是否再次尝试。\code{WATCH}观察的是Redis键的变化，外部数据库中的版本变化要通过应用协议反映到这里。

事务要区分排队错误和执行错误。假定事务先把\code{projection:P1}设为字符串，再对它执行列表弹出\code{LPOP}，最后更新版本键。第二条在执行时返回\code{WRONGTYPE}，第一条已经成功的字符串写入保留，第三条继续执行，\code{EXEC}的结果数组同时包含成功和错误。若排队阶段就存在参数错误，服务器记录这一错误，\code{EXEC}以\code{EXECABORT}拒绝执行队列。客户端分别处理整个队列未执行与队列执行后部分子命令出错这两种结果。

\CEighteenFigure{ch18-20-transaction}{事务排队失败与执行期类型错误具有不同出口}{fig:18-transaction-branches}

脚本把读取、比较和写入放进服务器端的一次调用，适合实现投影版本比较或同槽状态判断。Lua脚本可通过\code{EVAL}传入脚本文本、键数量、键与其他参数，并以返回值给出比较或写入结果。脚本运行期间，普通客户端命令不会插进它的步骤之间。应用先校验当前类型、输入和必要条件，再安排能够完成的状态修改。脚本写入一部分数据后发生运行错误，已经完成的修改保留。

长脚本占用命令执行路径，会延迟其他请求。超时处理允许的管理动作有各自条件：对于客户端直接运行的脚本，尚未写数据时可用相应终止命令停止，已经写过数据后则受到更严格的处理限制。超时本身不撤销已完成写入。事务和脚本的上述行为围绕Redis本地键空间展开，外部支付HTTP调用由工作者在既定C/T协议中完成~\cite{redis8102-source}。

\subsection{执行产生的修改怎样传播}

传播层是命令执行与后续保存、复制之间的衔接代码。它收集本次执行实际产生且需要传播的修改，一路编码成AOF记录，一路进入主节点向副本发送的复制流。AOF即追加文件，保存可供重启重放的数据变化；复制流把变化送到其他运行中的节点。下一节起分别展开文件与副本怎样接上这些修改。

客户端原始请求和传播记录可能不同。读取发现过期键时，传播可以包含删除；脚本根据读值选择分支后，传播反映实际执行的写命令。一个执行单元产生多项需要共同重放的修改时，Redis按相应规则用\code{MULTI}/\code{EXEC}包裹传播操作。加载AOF或处理复制流的节点据此连续应用这组操作，避免把同一组合拆成可随意插入其他命令的几段~\cite{redis8102-source}。

\CEighteenFigure{ch18-21-propagation}{本地执行收集实际修改，再向AOF与复制流传播}{fig:18-propagation-unit}

传播边界保持Redis数据的重放次序，应用的成功分支、冲突分支、原结果保存及未决责任则要先形成完整业务结构。例如，先用\code{SET NX}抢到一个S占位，再分别写入原输入、C、T和原结果，进程退出可能落在任意两次调用之间。传播层会忠实保存已经执行的那几步，恢复者面对的却可能是含义不完整的占位。把它设计成持久受理事实时，就同时保存绑定和恢复所需字段；把它设计成临时构造时，就在发布为可见受理以前完成构造及检查。

\subsection{陈旧回填与缺失项的更新资格}

查询缓存常用读取缺失后查数据库、回填缓存的路径。初态是数据库中P1的投影版本1、状态\code{paid}。读者A查到这一版本，回填动作被延迟。履约更新随后提交版本2、状态\code{fulfilled}，并删除缓存。A最后执行普通\code{SET}，把版本1重新放进已经删除的位置。此后命中缓存的读者会再次看到\code{paid}，直到下一次刷新或失效。

\CEighteenFigure{ch18-03}{数据库更新与失效之后，旧读者重新填入版本1}{fig:18-stale-fill}

Memcached是另一种常见的内存缓存服务器。本章采用1.6.45讨论它的缓存协议：\code{get}读取项，\code{set}存储项，\code{gets}另外返回一个不透明CAS值，\code{cas}写入时比较当前缓存项的CAS值。CAS即比较后交换，服务器只在条件匹配时更新。读者A取得CAS值c1以后，读者B替换了该项，A再用c1更新便得到\code{EXISTS}。该项被删除以后，A得到\code{NOT\_FOUND}。把失败后的\code{cas}降级为无条件\code{set}，又会允许旧版本填入缺失位置~\cite{memcached1645-protocol}。

针对缺失项，缓存系统可以另行发放回填租约。第一个miss读者取得与键和当前代次关联的token，后续读者等待或回源到受限路径。数据库更新后的失效操作同时撤销这份token，旧读者携带它回填时被拒绝。新的miss获得新token，读取版本2后才能回填。缓存节点重启时，token还要包含能够区分旧实例的成分，或者通过等效协议拒绝重启前的迟到回填。

\CEighteenFigure{ch18-04}{缺失项租约把回填资格与失效操作连接起来}{fig:18-cache-lease}

Facebook的2013年memcache设计用按键发放的租约处理陈旧回填和热点miss，把删除与租约失效联系起来，并控制发放速度。当前Memcached的meta协议也提供缺失时建立stub、返回\code{W}给一个回填者、返回\code{Z}表示已有获胜者的机制，还能结合CAS约束更新\footnote{stub是在缓存中为缺失项暂建的占位对象，用于关联当前回填资格。这里的\code{W}/\code{Z}是meta协议中的标记，实际命令还要按1.6.45的选项、CAS和删除语义组织。}。两者属于不同时间和实现中的具体协议，早期租约定义中的接口不能原样套到普通\code{get}/\code{set}上~\cite{memcache-facebook2013,memcached1645-protocol}。

租约的有效性还依赖失效动作何时到达。数据库已经提交版本2，负责删除缓存的进程却在发送删除前退出，那么原租约可能继续有效，缓存旧值也可能继续被读到。购买系统可把待失效事件与权威投影变更一起保存，再由恢复者持续投递，形成可恢复的失效责任。提交到失效完成之间的读取窗口，由查询接口明确处理。要求读到某次写入后状态的调用方，可以携带最低版本，读到较旧缓存时回源；原付款结果查询直接从原S的权威记录恢复。

\CEighteenFigure{ch18-22-invalidation}{持久失效责任跨过进程退出，最低版本约束处理传播窗口}{fig:18-invalidation-gap}

热点miss还涉及负载。一个回填者得到资格后，其余读者可以短暂等待，允许陈旧展示的页面也可以在后台刷新时读取明确标注的旧投影。资金和履约执行按其权威状态判断。回填失败、租约过期以及工作者退出以后，后续读者重新竞争资格，TTL界定这次缓存协调的寿命。

失效投递本身可以重复。某次删除已经执行，响应丢失，恢复者又删除一次，结果通常只是让后续读取重新回源。持续重复可能造成无谓回源，并反复撤销新token，失效事件因此按来源版本合并或收敛，避免较旧失效持续撤销后来合法取得的回填资格。

版本约束让页面读取具有明确的推进条件。调用方刚收到投影版本2，之后查询携带\code{min\_revision=2}，缓存命中版本1时，查询服务转回能够提供足够版本的来源。来源尚未应用到版本2时，可以等待到达或返回明确的进度状态。普通页面查询则按该接口允许的陈旧范围返回。版本值表达投影进度，原操作接口另用S保存不可覆盖的原结果。
\FloatBarrier
